% M1 report skeleton — ACM sigconf template.
% Compile on Overleaf (search "ACM sigconf" template, or upload this + references.bib) —
% the acmart class isn't installed locally. Word template is also available from ACM if
% the group prefers that instead; this file assumes LaTeX.

\documentclass[sigconf]{acmart}

% Suppress the ACM copyright block / DOI / ISBN stamp — this is a course report,
% not an actual ACM publication.
\settopmatter{printacmref=false}
\renewcommand\footnotetextcopyrightpermission[1]{}
\pagestyle{plain}
\acmConference[PRI 2026/27]{Information Processing and Retrieval}{2026}{Porto, Portugal}

\begin{document}

\title{Find Movies Like the One You're Thinking Of}
\subtitle{An Information Retrieval System Linking Movie Metadata and Wikipedia Plots}

% TODO: fill in real names/emails for all 4 group members.
\author{Lara}
\affiliation{%
  \institution{FEUP, U.Porto}
  \country{Portugal}
}

\author{P2 Name}
\affiliation{%
  \institution{FEUP, U.Porto}
  \country{Portugal}
}

\author{P3 Name}
\affiliation{%
  \institution{FEUP, U.Porto}
  \country{Portugal}
}

\author{P4 Name}
\affiliation{%
  \institution{FEUP, U.Porto}
  \country{Portugal}
}

\begin{abstract}
Existing movie search tools let users filter by title, genre, cast, or year, but fail
when a user remembers only a vague plot detail rather than any exact metadata. We
present a document collection and retrieval system that combines \textit{The Movie
Database (TMDB) Comprehensive Dataset} (Kaggle), providing structured movie metadata
and user reviews, with plot summaries scraped from Wikipedia and linked through
Wikidata. Each document describes a single movie through both structured fields and
substantial free text, enabling retrieval that goes beyond exact-field matching to
surface movies by narrative and thematic description. This report covers the data
sources, the collection and linking pipeline, dataset characterization, and a set of
prospective information needs motivating the retrieval work in later milestones.
\end{abstract}

\maketitle

\section{Topic}
% The differentiating angle: why movies despite it being a common topic.
Movies are a common dataset choice, so this project deliberately does not center on
metadata search (genre, year, cast) that any catalog already supports well. Instead,
the goal is retrieval driven by \emph{what a user remembers about a movie} — a plot
detail, a theme, a narrative trope, or what critics said about it — rather than what
they can already filter for. This motivates linking structured TMDB metadata with two
free-text sources: Wikipedia plot summaries (narrative content) and TMDB user reviews
(reception and opinion), so that information needs can be phrased the way people
actually half-remember movies, not the way a database schema represents them.

\section{Data Sources}

\subsection{The Movie Database (TMDB) Comprehensive Dataset}
Kaggle dataset by rishabhkumar2003, sourced from the TMDB API. Five CSV files:
\texttt{movies.csv} (9{,}771 rows — one per movie; title, overview, release date,
runtime, budget, revenue, rating, vote count, genres, language), \texttt{cast.csv}
(\textasciitilde150{,}000 rows), \texttt{crew.csv} (\textasciitilde63{,}600 rows),
\texttt{genres.csv} (a 20-row reference table, not a per-movie join table — genres are
already embedded in \texttt{movies.csv}), and \texttt{reviews.csv} (user reviews,
13{,}073 rows after parsing). No IMDb ID field exists in this dataset.
% TODO: confirm and state the dataset's license explicitly.

\subsection{Wikipedia}
Plot summaries for each movie, retrieved via the Wikipedia/MediaWiki API
(\texttt{action=parse}), linked to TMDB IDs through Wikidata (property P4947) rather
than by title matching, to avoid collisions with same-titled non-film articles.
Wikipedia text is licensed CC~BY-SA~4.0.
% TODO: confirm exact license wording to cite.

\section{Collection \& Preparation}
% TODO: insert pipeline diagram (UML-style) once the full pipeline is merged.
The collection pipeline has four stages, run by different group members and merged
into a single final collection:
\begin{enumerate}
  \item \textbf{Movies/reviews cleaning} — deduplicate movies, drop movies with no
    release date, clean review text (strip HTML/Markdown formatting), drop empty or
    very short (\textless20 words) reviews, drop duplicate reviews. Result:
    9{,}771 $\rightarrow$ 9{,}710 movies, 13{,}073 $\rightarrow$ 12{,}248 reviews.
  \item \textbf{Wikidata linking} — resolve each TMDB ID to an English Wikipedia
    article title via Wikidata property P4947. 7{,}016 / 9{,}710 movies (72.3\%)
    linked successfully.
  \item \textbf{Wikipedia fetching \& parsing} — fetch each linked article and extract
    its Plot section.
    % TODO: numbers once the fetch runs on the full linked set with the corrected titles.
  \item \textbf{Final join} — merge cleaned movies, reviews, cast/crew, Wikidata
    links, and Wikipedia plots into one document per movie.
    % TODO: not yet complete.
\end{enumerate}

% TODO: conceptual data model (document = movie, has reviews, has cast/crew, has a
% linked Wikipedia plot, etc.)

\section{Characterization}
% Reviews characterization (P1's part) is complete; the rest needs the final
% joined collection.

\subsection{Reviews}
Of 9{,}710 cleaned movies, 6{,}013 (62\%) have zero reviews after cleaning; review
coverage is sparse and concentrated on popular titles. Movies with at least one review
average 1.26 reviews (median 0 across all movies, since most have none); reviews
average 239 words (median 181). Review length shows essentially no correlation with a
movie's rating (Pearson $r = 0.115$, $n=3{,}697$). The most frequent words across all
reviews, after stopword removal, are dominated by generic film-review vocabulary
(\textit{story, character, cast, plot, scenes}) rather than genre-distinctive terms —
see \texttt{data/interim/movies\_reviews\_report.md} for full tables.

% TODO: collection-wide stats (doc count, vocabulary size, frequent terms),
% structured-field stats (genre/year/language/rating distributions),
% multi-variable analysis, Wikipedia-section stats, NER — all pending the final join.

\section{Prospective Search Tasks}
Information needs are deliberately phrased around themes, situations, and reception
rather than metadata filters, since the latter is already well-served by existing
movie catalogs.

\begin{enumerate}
  \item A movie whose reviewers describe the comedy as clever, witty, or
    intellectually sharp, rather than slapstick or lowbrow.
  \item A romantic movie where reviewers are divided — some call it predictable or
    clichéd, others find it genuinely heartwarming.
  \item A science-fiction survival film about astronauts stranded away from Earth who
    need to survive or get home.
  \item A crime or thriller film in which an FBI agent investigates serial killings.
\end{enumerate}
% TODO: 5 more needed (2 from P3, 3 from P4, including one combining a text query
% with a structured filter) — see docs/information_needs.md for the living list.

\section{References}
% See references.bib — TMDB/Kaggle dataset, Wikipedia, Wikidata, MediaWiki API.
\bibliographystyle{ACM-Reference-Format}
\bibliography{references}

\end{document}
